% !TeX program = xelatex
\documentclass[11pt,a4paper]{article}

\usepackage[a4paper,margin=24mm]{geometry}
\usepackage{kotex}
\usepackage{amsmath,amssymb,mathtools,bm}
\usepackage{booktabs,longtable,array}
\usepackage{enumitem}
\usepackage{xcolor}
\usepackage{hyperref}

\hypersetup{
  colorlinks=true,
  linkcolor=blue!55!black,
  urlcolor=blue!55!black,
  pdftitle={BARI2D 수학적 정식화},
  pdfauthor={BARI2D Project}
}

\newcommand{\R}{\mathbb{R}}
\newcommand{\N}{\mathbb{N}}
\newcommand{\I}{\mathbb{I}}
\newcommand{\Ftarget}{F_{\mathrm{target}}}
\newcommand{\Fcap}{F_{\mathrm{cap}}}
\newcommand{\Pspan}{P_{\mathrm{span}}}
\newcommand{\Ispan}{I_{\mathrm{span}}}
\newcommand{\Gcontact}{\mathcal{G}_{\mathrm{contact}}}
\newcommand{\norm}[1]{\left\lVert #1\right\rVert}
\newcommand{\clip}{\operatorname{clip}}
\newcommand{\concat}{\operatorname{concat}}
\newcommand{\sigmoid}{\operatorname{sigmoid}}
\newcommand{\softmax}{\operatorname{softmax}}
\newcommand{\MSE}{\operatorname{MSE}}

\title{\textbf{BARI2D 수학적 정식화}\\[4pt]
\large 분산 다중 로봇 교량 건설의 상태, 동역학, 구조 용량 및 MAPPO 목적}
\author{BARI2D Project}
\date{\today}

\begin{document}
\maketitle

\begin{abstract}
이 문서는 BARI2D 구현의 수학적 구조를 정리한다. 환경 좌표계, 로봇 상태와 이산 동역학, 국소 감각, 방향성 사각형 접촉, 앵커 파손, 기계 접촉 그래프, 제방 간 연결도, 그래프 기반 하중 용량, 기능 보상, 생물 영감형 공유 정책, 중앙 비평가와 순환형 MAPPO 목적을 정의한다. 본 정식화의 핵심 제약은 분산 행위자가 전역 지도, 전역 위치, 접촉 그래프와 현재 구조 용량을 관측하지 않는다는 것이다. 특정 교량 형상이나 트러스 토폴로지는 목적 함수에 포함하지 않는다.
\end{abstract}

\tableofcontents

\newpage

\section{문제 정의}

\subsection{다중 로봇 협동 문제}

로봇 집합을
\begin{equation}
  \mathcal{R}=\{1,\ldots,N\}, \qquad N\approx 20
\end{equation}
으로 둔다. 모든 로봇은 동일한 물리 파라미터와 정책 파라미터 \(\theta\)를 공유한다. 에피소드마다 목표 하중은
\begin{equation}
  \Ftarget \sim p_F(F), \qquad
  F\in[F_{\min},F_{\max}]
\end{equation}
에서 샘플링된다.

로봇 \(i\)의 분산 정책은
\begin{equation}
  a_i^t \sim
  \pi_\theta\!\left(
    a_i^t\mid o_i^t,h_i^t,\Ftarget
  \right)
  \label{eq:decentralized-policy}
\end{equation}
로 정의한다. \(o_i^t\)는 국소 관측, \(h_i^t\)는 순환 은닉 상태다. 실행 중 중앙 제어기는 존재하지 않는다.

\subsection{최적화 목표}

기능적 목표는 다음의 다목적 최적화로 해석할 수 있다.
\begin{equation}
  \max_\theta\;
  \mathbb{E}_{\tau\sim\pi_\theta}
  \left[
    w_s S(\tau)
    -w_T T(\tau)
    -w_E E(\tau)
    -w_R N_{\mathrm{used}}(\tau)
  \right],
  \label{eq:global-objective}
\end{equation}
여기서 성공 지표는
\begin{equation}
  S(\tau)=
  \I\!\left[
    \Ispan(\tau)=1
    \;\land\;
    \Fcap(\tau)\ge\Ftarget
  \right]
  \label{eq:success}
\end{equation}
이다. 목적식에는 특정 형상, 각도 또는 토폴로지 항이 없다.

\section{필드와 간극 기하}

\subsection{전역 영역}

시뮬레이션 영역은
\begin{equation}
  \Omega=[0,L_x]\times[0,L_y]\subset\R^2
\end{equation}
다. 간극 중심을 \(\bm{c}\in\R^2\), 간극을 가로지르는 단위 법선을
\begin{equation}
  \bm{n}=
  \begin{bmatrix}
    \cos\psi\\ \sin\psi
  \end{bmatrix}
\end{equation}
로 둔다. 접선은
\begin{equation}
  \bm{t}=
  \begin{bmatrix}
    -\sin\psi\\ \cos\psi
  \end{bmatrix}
\end{equation}
다.

점 \(\bm{p}\in\Omega\)의 간극 좌표는
\begin{align}
  s(\bm{p}) &= (\bm{p}-\bm{c})^\top\bm{n},\\
  q(\bm{p}) &= (\bm{p}-\bm{c})^\top\bm{t}.
\end{align}

\subsection{불규칙 경계}

간극 폭을 \(w\), 경계 섭동 진폭을 \(A\), 공간 주파수를 \(\omega\), 위상을 \(\phi\)라 하면 양쪽 경계는
\begin{align}
  b_L(q) &= -\frac{w}{2}+A\sin(\omega q+\phi),\\
  b_R(q) &=  \frac{w}{2}+A\sin(\omega q+\phi)
\end{align}
로 정의한다. 제방 및 간극 집합은
\begin{align}
  \Omega_L &= \{\bm{p}\in\Omega:s(\bm{p})\le b_L(q(\bm{p}))\},\\
  \Omega_R &= \{\bm{p}\in\Omega:s(\bm{p})\ge b_R(q(\bm{p}))\},\\
  \Omega_G &= \Omega\setminus(\Omega_L\cup\Omega_R)
\end{align}
이다.

점의 정규화 간극 진행도는
\begin{equation}
  \rho(\bm{p})=
  \clip\!\left(
    \frac{s(\bm{p})-b_L(q(\bm{p}))}
         {b_R(q(\bm{p}))-b_L(q(\bm{p}))},
    0,1
  \right)
  \label{eq:point-progress}
\end{equation}
로 둔다.

\section{로봇 상태와 동역학}

\subsection{상태 벡터}

로봇 \(i\)의 물리 및 내부 상태를
\begin{equation}
  x_i^t=
  \left(
    \bm{p}_i^t,\theta_i^t,v_i^t,\delta_i^t,
    \ell_i^t,b_i^t,\varepsilon_i^t,
    a_i^{t-1},k_i^t,f_i^t
  \right)
\end{equation}
로 둔다. 각 성분은 다음과 같다.

\begin{center}
\begin{tabular}{cl}
\toprule
기호 & 의미\\
\midrule
\(\bm{p}_i=(x_i,y_i)\) & 중심 위치\\
\(\theta_i\) & 방향각\\
\(v_i\) & 전진/후진 속도\\
\(\delta_i\) & 조향각\\
\(\ell_i\in\{0,1\}\) & 헤드 리프트 상태\\
\(b_i\in\{0,1\}\) & 앵커 상태\\
\(\varepsilon_i\ge0\) & 국소 변형률/하중 추정\\
\(a_i^{t-1}\) & 이전 행동\\
\(k_i\in\{0,\ldots,K\},\;K=10\) & 2.5D 높이 레이어\\
\(f_i\in\{0,1\}\) & 낙하 상태\\
\bottomrule
\end{tabular}
\end{center}

\subsection{평면 운동학}

로봇 진행 단위벡터는
\begin{equation}
  \bm{d}(\theta_i)=
  \begin{bmatrix}
    \cos\theta_i\\ \sin\theta_i
  \end{bmatrix}
\end{equation}
다. 이산 이동 명령에서 속도 명령 \(u_{v,i}^t\in[-1,1]\), 조향 명령 \(u_{\delta,i}^t\in[-1,1]\)을 얻는다. 구동 잡음은
\begin{equation}
  \widetilde{u}_{v,i}^t
  =\clip(u_{v,i}^t+\eta_{v,i}^t,-1,1),
  \qquad
  \eta_{v,i}^t\sim\mathcal{N}(0,\sigma_a^2)
\end{equation}
로 표현한다. 조향도 같은 방식으로 잡음을 추가한다.

운동학 업데이트는
\begin{align}
  v_i^{t+1}
    &=v_{\max}\widetilde{u}_{v,i}^t,\\
  \delta_i^{t+1}
    &=\delta_{\max}\widetilde{u}_{\delta,i}^t,\\
  \theta_i^{t+1}
    &=\operatorname{wrap}\!\left(
      \theta_i^t+
      \omega_{\max}\widetilde{u}_{\delta,i}^t\Delta t
    \right),\\
  \bm{p}_i^{t+1}
    &=\bm{p}_i^t+
      v_i^{t+1}\bm{d}(\theta_i^{t+1})\Delta t
  \label{eq:kinematics}
\end{align}
다.

이동 에너지 프록시는
\begin{equation}
  \Delta E_i^t=
  |v_i^{t+1}|\Delta t+
  c_\delta|\delta_i^{t+1}|\Delta t
\end{equation}
로 누적한다. 등반 시에는
\begin{equation}
  \Delta E_{i,\mathrm{climb}}=m_i g_{\mathrm{proxy}}h_c
\end{equation}
에 대응하는 질량--높이 프록시를 추가한다.

\subsection{이산 등반}

로봇 \(j\)가 로봇 \(i\)의 진행 방향 앞에 있고
\begin{equation}
  k_j^t=k_i^t,
  \qquad
  \norm{\bm{p}_j-\bm{p}_i}\le r_c,
  \qquad
  (\bm{p}_j-\bm{p}_i)^\top\bm{d}(\theta_i)>0
\end{equation}
이면 \(j\)는 \(i\)의 등반 지지가 될 수 있다. 등반 업데이트는
\begin{equation}
  k_i^{t+1}=
  \min(k_j^t+1,K)
\end{equation}
다. 이는 연속 \(z\)축 강체 운동 대신 사용하는 2.5D 근사다.

등반 뒤 지지체를 지나 전진한 로봇은 별도 하강 행동 없이 자동으로 내려온다.
\(\mathcal{B}_i^t\)를 로봇 \(i\)의 평면 방향성 본체 영역이라 하고,
\(k_i^t>0\)인 로봇과 양의 면적으로 겹치는 하위 레이어 지지 집합을
\begin{equation}
  \mathcal{S}_{i,\downarrow}^t=
  \left\{
    j:
    k_j^t<k_i^t,\;
    \operatorname{area}(\mathcal{B}_i^t\cap\mathcal{B}_j^t)>0,\;
    f_j^t=0
  \right\}
\end{equation}
라 하자. 모멘트나 앵커 여부를 사용하지 않고 레이어를
\begin{equation}
  k_i^{t+1}=
  \begin{cases}
    \min\left(k_i^t,1+\max_{j\in\mathcal{S}_{i,\downarrow}^t} k_j^t\right),
      & \mathcal{S}_{i,\downarrow}^t\ne\varnothing,\\
    0, & \mathcal{S}_{i,\downarrow}^t=\varnothing
  \end{cases}
\end{equation}
로 갱신한다. 모든 로봇의 레이어가 더 이상 변하지 않을 때까지 동시에 반복해
하부 지지체의 연쇄 낙하도 같은 step 안에 반영한다.

\section{국소 센서와 관측}

\subsection{IR 광선}

로봇 \(i\)의 \(r\)번째 센서 상대각을 \(\alpha_r\)라 하면 광선은
\begin{equation}
  \bm{r}_{i,r}(\lambda)=
  \bm{p}_i+
  \lambda\bm{d}(\theta_i+\alpha_r),
  \qquad
  \lambda\in[0,D_{\max}]
\end{equation}
다. 검출 집합을 다른 로봇, 장애물, 필드 경계와 지지면 전이의 합집합 \(\mathcal{O}_i^t\)라 두면
\begin{equation}
  d_{i,r}^t=
  \min\left\{
    \lambda\in[0,D_{\max}]
    :\bm{r}_{i,r}(\lambda)\in\mathcal{O}_i^t
  \right\}.
\end{equation}
검출이 없으면 \(d_{i,r}^t=D_{\max}\)다. 정책 입력은
\begin{equation}
  \bar d_{i,r}^t=
  \clip\!\left(
    \frac{d_{i,r}^t+\eta_{i,r}^t}{D_{\max}},
    0,1
  \right)
\end{equation}
이며 \(\eta_{i,r}^t\sim\mathcal{N}(0,\sigma_s^2)\)다.

기본 수평 센서 집합은
\begin{equation}
  \mathcal{A}=
  \left\{
    0,\pi,\frac{\pi}{2},-\frac{\pi}{2}
  \right\}
\end{equation}
로, 각각 전방, 후방, 좌측, 우측을 뜻한다. 여기에 하향 IR
\(\bar d_{i,\downarrow}^t\)를 추가한다. 하향 센서는 로봇 중심 아래 제방 또는 반경
\(r_c\) 안의 더 낮은 레이어 로봇까지의 수직 거리를 측정한다.
\begin{equation}
  d_{i,\downarrow}^t=
  \begin{cases}
    k_i^t h_c, & \text{중심 아래가 제방인 경우},\\
    \min\limits_{j:\;k_j^t<k_i^t,\;
      \norm{\bm{p}_j^t-\bm{p}_i^t}\le r_c}
      (k_i^t-k_j^t)h_c, & \text{아래 로봇이 있는 경우},\\
    D_{\max}, & \text{아래 표면이 없는 경우}.
  \end{cases}
\end{equation}
따라서 \(\bar d_{i,\downarrow}^t=1\)은 절벽 또는 지지 상실의 국소 신호다.

\subsection{변형률}

로봇 \(i\)와 연결된 접촉 및 앵커 간선 집합을
\(\mathcal{E}_i^t\)라 두면 국소 하중은
\begin{equation}
  \varepsilon_i^t=
  \sum_{e\in\mathcal{E}_i^t}|F_e^t|
\end{equation}
로 근사한다. 정규화 입력은
\begin{equation}
  \bar\varepsilon_i^t=
  \clip\!\left(
    \frac{\varepsilon_i^t+\eta_{\varepsilon,i}^t}
         {F_{\mathrm{scale}}},
    0,\varepsilon_{\max}
  \right)
\end{equation}
다.

\subsection{시간 이력과 관측 벡터}

길이 \(H\)의 이력은
\begin{align}
  D_i^t &=
  [\bar{\bm d}_i^{t-H+1},\ldots,\bar{\bm d}_i^t],\\
  E_i^t &=
  [\bar\varepsilon_i^{t-H+1},\ldots,\bar\varepsilon_i^t],\\
  A_i^t &=
  [\bar a_i^{t-H},\ldots,\bar a_i^{t-1}]
\end{align}
다. 내부 상태를
\begin{equation}
  s_{i,\mathrm{int}}^t=
  \left[
    b_i^t,\ell_i^t,
    \frac{v_i^t}{v_{\max}},
    \frac{\delta_i^t}{\delta_{\max}},
    \frac{k_i^t}{K},
    f_i^t
  \right]
\end{equation}
로 두면 최종 국소 관측은
\begin{equation}
  o_i^t=
  \concat\!\left(
    D_i^t,E_i^t,A_i^t,
    s_{i,\mathrm{int}}^t,
    \frac{\Ftarget}{F_{\mathrm{norm}}},
    z_i
  \right).
  \label{eq:local-observation}
\end{equation}
여기서
\begin{equation}
  z_i\sim\mathcal{N}(0,\sigma_z^2I)
\end{equation}
는 에피소드 시작 시 한 번 샘플링하며 기본값은 \(\sigma_z=0\)이다.

\section{접촉 역학과 앵커}

\subsection{방향성 사각형 접촉}

로봇 \(i\)의 방향성 사각형을 \(\mathcal{B}_i\)라 하자. 두 사각형의 모서리 법선에서 얻은 분리축 집합을 \(\mathcal{A}_{ij}\)라 두면
\begin{equation}
  \mathcal{B}_i\cap\mathcal{B}_j\ne\varnothing
  \iff
  \forall\bm{a}\in\mathcal{A}_{ij},\;
  \operatorname{proj}_{\bm a}(\mathcal{B}_i)
  \cap
  \operatorname{proj}_{\bm a}(\mathcal{B}_j)
  \ne\varnothing.
\end{equation}

축 \(\bm a\)의 겹침 길이를 \(\delta_{\bm a}\)라 하면 최소 침투와 접촉 법선은
\begin{equation}
  \delta_{ij}=\min_{\bm a\in\mathcal{A}_{ij}}\delta_{\bm a},
  \qquad
  \bm n_{ij}=
  \arg\min_{\bm a\in\mathcal{A}_{ij}}\delta_{\bm a}.
\end{equation}
법선 접촉력 프록시는
\begin{equation}
  F_{n,ij}=
  \max\!\left(
    0,\,
    k_n\delta_{ij}
    +c_n|v_i-v_j|
  \right)
  \label{eq:normal-force}
\end{equation}
다. 위치 보정 후 속도는 마찰 계수 \(\mu\)로 감쇠한다.

\subsection{앵커 힘과 파손}

로봇 \(i\)와 상대 \(j\) 사이 앵커의 현재 상대벡터를
\begin{equation}
  \bm r_{ij}=\bm p_j-\bm p_i,
  \qquad
  d_{ij}=\norm{\bm r_{ij}}
\end{equation}
로 둔다. 생성 시 기준 거리를 \(d_{ij}^0\)라 하면 축력은
\begin{equation}
  F_{\mathrm{axial}}=
  k_a(d_{ij}-d_{ij}^0)
\end{equation}
이다. \(F_{\mathrm{axial}}>0\)이면 인장, 음수면 압축이다.

앵커 방향 단위벡터
\(\widehat{\bm r}_{ij}=\bm r_{ij}/d_{ij}\)를 이용한 전단 프록시는
\begin{equation}
  F_{\mathrm{shear}}=
  c_s\left|
    d_x(\theta_i)\widehat r_{ij,y}
    -d_y(\theta_i)\widehat r_{ij,x}
  \right|
\end{equation}
다. 선택적 회전 항은
\begin{equation}
  M_{\mathrm{rot}}=k_{\mathrm{rot}}|\Delta\theta_{ij}|
\end{equation}
이다.

앵커 파손 지표는
\begin{equation}
  \chi_{ij}^{\mathrm{fail}}=
  \I\!\left[
    F_{\mathrm{axial}}>F_T
    \;\lor\;
    -F_{\mathrm{axial}}>F_C
    \;\lor\;
    F_{\mathrm{shear}}>F_S
  \right]
  \label{eq:anchor-failure}
\end{equation}
다. 확률 파손을 사용할 경우 베르누이 항을 추가한다.

\section{기계 접촉 그래프}

\subsection{그래프 정의}

시간 \(t\)의 기계 그래프는
\begin{equation}
  \Gcontact^t=(\mathcal{V},\mathcal{E}^t)
\end{equation}
이며
\begin{equation}
  \mathcal{V}=
  \mathcal{R}\cup\{L,R\}
\end{equation}
다. \(L,R\)은 왼쪽 및 오른쪽 제방 노드다.

각 간선 \(e=(u,v)\)는
\begin{equation}
  \xi_e=
  (\tau_e,c_e,k_e,F_e,\Delta\bm p_e,\Delta\theta_e)
\end{equation}
를 갖는다. \(\tau_e\)는 접촉, 앵커 또는 제방 접촉 유형이고 \(c_e\)는 용량, \(k_e\)는 강성이다.

\subsection{횡단 판정}

\(L\)에서 도달 가능한 노드 집합을
\begin{equation}
  \mathcal{C}_L^t=
  \operatorname{Reachable}(\Gcontact^t,L)
\end{equation}
라 하면
\begin{equation}
  \Ispan^t=
  \I[R\in\mathcal{C}_L^t]
  \label{eq:span-indicator}
\end{equation}
다.

연속 진행도는
\begin{equation}
  \Pspan^t=
  \max_{i\in\mathcal{C}_L^t\cap\mathcal{R}}
  \rho(\bm p_i^t),
  \label{eq:span-progress}
\end{equation}
이며 연결 로봇이 없으면 0이다.

\subsection{지지 상실}

양쪽 제방에 연결된 로봇 집합을
\begin{equation}
  \mathcal{S}^t=
  \operatorname{Reachable}(\Gcontact^t,L)
  \cup
  \operatorname{Reachable}(\Gcontact^t,R)
\end{equation}
로 둔다. 로봇 중심이 \(\Omega_G\)에 있고 몸체가 제방과 접촉하지 않으며
\(i\notin\mathcal{S}^t\)이면 낙하 상태로 전환한다.

\section{하중 용량}

\subsection{간선 용량 최대 유량}

그래프의 비음수 유량 \(f_{uv}\)에 대해 용량 제약은
\begin{equation}
  0\le f_{uv}\le c_{uv}
\end{equation}
이고 중간 노드의 유량 보존은
\begin{equation}
  \sum_{u:(u,v)\in\mathcal{E}}f_{uv}
  =
  \sum_{w:(v,w)\in\mathcal{E}}f_{vw}
\end{equation}
다.

중앙 하중 노드 \(m\)에서 각 제방으로 전달 가능한 최대 유량을
\begin{align}
  \Phi_L(m)&=\operatorname{MaxFlow}(m,L\mid R\ \text{차단}),\\
  \Phi_R(m)&=\operatorname{MaxFlow}(m,R\mid L\ \text{차단})
\end{align}
로 정의한다. 한 연결 구성요소의 양측 지지 용량은
\begin{equation}
  F_{\mathrm{flow}}^{(c)}
  =
  2\min\{\Phi_L(m_c),\Phi_R(m_c)\}
  \label{eq:component-flow}
\end{equation}
다. \(m_c\)는 해당 구성요소에서 \(|\rho(\bm p_i)-1/2|\)가 최소인 로봇이다.

균등 분산 하중에서는 모든 독립 횡단 구성요소 \(\mathcal{C}_{\mathrm{span}}\)의 용량을 합산한다.
\begin{equation}
  F_{\mathrm{flow}}^{\mathrm{uniform}}
  =
  \sum_{c\in\mathcal{C}_{\mathrm{span}}}
  F_{\mathrm{flow}}^{(c)}.
  \label{eq:uniform-capacity}
\end{equation}

\subsection{경로 강성과 변위 제한}

경로 \(p\)에 직렬로 놓인 간선들의 등가 강성은
\begin{equation}
  k_p=
  \left(
    \sum_{e\in p}\frac{1}{k_e}
  \right)^{-1}
  \label{eq:path-stiffness}
\end{equation}
이다. 독립 병렬 경로의 등가 강성은
\begin{equation}
  k_{\mathrm{eq}}=\sum_p k_p
\end{equation}
다. 최대 허용 변위가 \(d_{\max}\)이면 변위 기반 용량은
\begin{equation}
  F_{\mathrm{disp}}=k_{\mathrm{eq}}d_{\max}
\end{equation}
다.

빠른 평가 용량은
\begin{equation}
  \widehat{\Fcap}
  =
  \begin{cases}
    \min(F_{\mathrm{flow}},F_{\mathrm{disp}}),
      & \Ispan=1,\\
    0, & \Ispan=0
  \end{cases}
  \label{eq:fast-capacity}
\end{equation}
로 정의한다.

\subsection{증분 종단 시험}

하중 간격을 \(\Delta F\)라 하고
\begin{equation}
  F^{(n)}=n\Delta F
\end{equation}
로 하중을 증가시킨다. 다음 실패 조건을 검사한다.

\begin{align}
  \mathcal{H}_{\mathrm{anchor}}
    &: F^{(n)}\text{가 앵커 간선 용량을 초과},\\
  \mathcal{H}_{\mathrm{contact}}
    &: F^{(n)}\text{가 접촉 간선 용량을 초과},\\
  \mathcal{H}_{\mathrm{slip}}
    &: F^{(n)}\text{가 제방 접촉 용량을 초과},\\
  \mathcal{H}_{\mathrm{disp}}
    &: F^{(n)}/k_{\mathrm{eq}}>d_{\max},\\
  \mathcal{H}_{\mathrm{disconnect}}
    &: \Ispan=0.
\end{align}

첫 실패 직전의 지속 가능한 하중을
\begin{equation}
  \Fcap=
  \max\left\{
    F^{(n)}:
    \neg\mathcal{H}_{\mathrm{anchor}}
    \land
    \neg\mathcal{H}_{\mathrm{contact}}
    \land
    \neg\mathcal{H}_{\mathrm{slip}}
    \land
    \neg\mathcal{H}_{\mathrm{disp}}
    \land
    \neg\mathcal{H}_{\mathrm{disconnect}}
  \right\}
  \label{eq:terminal-capacity}
\end{equation}
로 둔다.

\section{기능 보상}

\subsection{연결 진행 보상}

\begin{equation}
  r_{\mathrm{span}}^t=
  \alpha\left(
    \Pspan^{t+1}-\Pspan^t
  \right).
  \label{eq:span-reward}
\end{equation}

\subsection{기계 성능 보상}

목표 대비 포화 품질을
\begin{equation}
  Q_t=
  \min\left(
    \frac{\widehat{\Fcap}^{\,t}}{\Ftarget},
    1
  \right)
\end{equation}
로 두면
\begin{equation}
  r_{\mathrm{mech}}^t=
  \beta(Q_{t+1}-Q_t).
  \label{eq:mechanical-reward}
\end{equation}

\subsection{효율과 실패}

\begin{equation}
  r_{\mathrm{eff}}^t=
  -\lambda_t
  -\lambda_E\Delta E^t
  -\lambda_R\Delta N_{\mathrm{used}}^t
  -\lambda_A\Delta N_{\mathrm{anchor}}^t
  -\lambda_C\Delta N_{\mathrm{fallen}}^t.
  \label{eq:efficiency-reward}
\end{equation}

종단 성공 보상은
\begin{equation}
  r_{\mathrm{success}}^t=
  R_S
  \I\!\left[
    \Ispan^t=1
    \land
    \Fcap^t\ge\Ftarget
  \right].
  \label{eq:success-reward}
\end{equation}

전체 팀 보상은
\begin{equation}
  r_t=
  r_{\mathrm{span}}^t+
  r_{\mathrm{mech}}^t+
  r_{\mathrm{eff}}^t+
  r_{\mathrm{success}}^t.
  \label{eq:team-reward}
\end{equation}

\section{모듈형 공유 행위자}

\subsection{교통 가지}

IR와 행동 이력을 시간축 입력으로 묶는다.
\begin{equation}
  X_{i,\mathrm{traffic}}^t
  =
  [D_i^t,A_i^t].
\end{equation}
두 개의 소형 1차원 시간 합성곱을 사용해
\begin{equation}
  z_{i,\mathrm{traffic}}^t
  =
  \operatorname{MeanPool}
  \left(
    \tanh\left(
      W_{t,2}*
      \tanh(W_{t,1}*X_{i,\mathrm{traffic}}^t+b_{t,1})
      +b_{t,2}
    \right)
  \right)
  \label{eq:traffic-encoder}
\end{equation}
을 만든다.

\subsection{연결성 가지}

\begin{equation}
  X_{i,\mathrm{conn}}^t
  =
  [D_i^t,E_i^t,A_i^t,\bar v_i^t]
\end{equation}
를 시간 GRU에 입력한다.
\begin{equation}
  z_{i,\mathrm{conn}}^t
  =
  \operatorname{GRU}_{\mathrm{conn}}
  (X_{i,\mathrm{conn}}^t).
  \label{eq:connectivity-encoder}
\end{equation}

\subsection{기계 가지와 목표 가지}

\begin{align}
  z_{i,\mathrm{mech}}^t
    &=
    \operatorname{TCN}_{\mathrm{mech}}(E_i^t,A_i^t),\\
  z_{\mathrm{goal}}
    &=
    \tanh\left(
      W_g\frac{\Ftarget}{F_{\mathrm{norm}}}+b_g
    \right).
\end{align}

\subsection{FiLM 조건화}

\begin{align}
  [\gamma_m,\beta_m]
    &=W_m^{\mathrm{film}}z_{\mathrm{goal}}+b_m^{\mathrm{film}},\\
  [\gamma_t,\beta_t]
    &=W_t^{\mathrm{film}}z_{\mathrm{goal}}+b_t^{\mathrm{film}},\\
  z'_{i,\mathrm{mech}}
    &=\gamma_m\odot z_{i,\mathrm{mech}}+\beta_m,\\
  z'_{i,\mathrm{traffic}}
    &=\gamma_t\odot z_{i,\mathrm{traffic}}+\beta_t.
  \label{eq:film}
\end{align}

\subsection{융합과 순환 상태}

\begin{equation}
  z_i^t=
  \concat\left(
    z'_{i,\mathrm{traffic}},
    z_{i,\mathrm{conn}},
    z'_{i,\mathrm{mech}},
    z_{\mathrm{goal}},
    s_{i,\mathrm{int}}^t,
    z_i
  \right).
\end{equation}
융합 특징은
\begin{equation}
  y_i^t=\tanh(W_fz_i^t+b_f)
\end{equation}
이고 순환 코어는 표준 GRU다.
\begin{align}
  r_i^t &=
  \sigmoid(W_ry_i^t+U_rh_i^{t-1}+b_r),\\
  u_i^t &=
  \sigmoid(W_uy_i^t+U_uh_i^{t-1}+b_u),\\
  \widetilde h_i^t &=
  \tanh(W_hy_i^t+U_h(r_i^t\odot h_i^{t-1})+b_h),\\
  h_i^t &=
  (1-u_i^t)\odot h_i^{t-1}
  +u_i^t\odot\widetilde h_i^t.
  \label{eq:actor-gru}
\end{align}

행동 로짓과 확률은
\begin{equation}
  \ell_i^t=W_\pi h_i^t+b_\pi,
  \qquad
  \pi_\theta(a_i^t\mid o_i^t,h_i^{t-1},\Ftarget)
  =
  \softmax(\ell_i^t+M_i^t)
  \label{eq:policy-head}
\end{equation}
다. \(M_{i,a}^t=0\)이면 가능한 행동이고,
\(M_{i,a}^t=-\infty\)이면 물리적으로 불가능한 행동이다.

\section{보조 표현 학습}

시뮬레이터 훈련 표적을
\begin{align}
  y_{i,\mathrm{conn}}^t
    &=\text{국소 로봇 접촉 이웃 수},\\
  y_{i,\mathrm{traffic}}^t
    &=\text{이동 중인 접촉 이웃 수},\\
  y_{i,\mathrm{persist}}^t
    &=\text{이전 스텝부터 유지된 접촉 수},\\
  y_{i,\mathrm{force}}^t
    &=\bar\varepsilon_i^t-\bar\varepsilon_i^{t-1}
\end{align}
로 둔다. 각 잠재 가지의 선형 헤드가 이를 예측한다.

\begin{align}
  \mathcal{L}_{\mathrm{conn}}
    &=\MSE(\widehat y_{\mathrm{conn}},y_{\mathrm{conn}}),\\
  \mathcal{L}_{\mathrm{traffic}}
    &=\MSE(\widehat y_{\mathrm{traffic}},y_{\mathrm{traffic}}),\\
  \mathcal{L}_{\mathrm{persist}}
    &=\MSE(\widehat y_{\mathrm{persist}},y_{\mathrm{persist}}),\\
  \mathcal{L}_{\mathrm{force}}
    &=\MSE(\widehat y_{\mathrm{force}},y_{\mathrm{force}}).
\end{align}

전체 보조 손실은
\begin{equation}
  \mathcal{L}_{\mathrm{aux}}=
  \lambda_{\mathrm{conn}}\mathcal{L}_{\mathrm{conn}}
  +\lambda_{\mathrm{traffic}}\mathcal{L}_{\mathrm{traffic}}
  +\lambda_{\mathrm{persist}}\mathcal{L}_{\mathrm{persist}}
  +\lambda_{\mathrm{force}}\mathcal{L}_{\mathrm{force}}.
  \label{eq:auxiliary-loss}
\end{equation}
이 표적은 추론 입력이 아니다.

\section{중앙 비평가}

\subsection{평탄화 비평가}

중앙 상태는
\begin{equation}
  s_t=
  \concat\left(
    \{g_i^t\}_{i=1}^N,
    w,\sin\psi,\cos\psi,
    \frac{\Ftarget}{F_{\mathrm{norm}}},
    \Pspan^t,
    \frac{\widehat{\Fcap}^{\,t}}{F_{\mathrm{norm}}},
    \Ispan^t
  \right)
\end{equation}
다. \(g_i^t\)는 위치, 방향, 변형률, 앵커, 속도, 레이어와 낙하를 포함한다. MLP 비평가는
\begin{equation}
  V_\varphi(s_t)=
  W_{v,3}\tanh\left(
    W_{v,2}\tanh(W_{v,1}s_t+b_{v,1})+b_{v,2}
  \right)+b_{v,3}
  \label{eq:mlp-critic}
\end{equation}
로 값을 예측한다.

\subsection{그래프 비평가}

로봇 노드 특징을 \(x_i^{(0)}\), 간선 특징을 \(\xi_{ij}\)라 하자. 메시지 전달층은
\begin{align}
  m_i^{(\ell)}
    &=
    \frac{1}{\max(1,|\mathcal{N}(i)|)}
    \sum_{j\in\mathcal{N}(i)}
    \left[
      W_n^{(\ell)}x_j^{(\ell)}
      \odot
      \sigmoid(W_e^{(\ell)}\xi_{ij}+b_e^{(\ell)})
    \right],\\
  x_i^{(\ell+1)}
    &=
    \tanh\left(
      W_s^{(\ell)}x_i^{(\ell)}
      +m_i^{(\ell)}
      +b_s^{(\ell)}
    \right).
  \label{eq:message-passing}
\end{align}

그래프 풀링과 가치 출력은
\begin{align}
  \bar x&=\frac{1}{N}\sum_{i=1}^Nx_i^{(L)},\\
  V_\varphi(\Gcontact^t)
    &=W_{g,2}\tanh(W_{g,1}\bar x+b_{g,1})+b_{g,2}.
  \label{eq:graph-value}
\end{align}

두 비평가 모두 훈련 전용이며 행위자의 분산 실행 입력을 바꾸지 않는다.

\section{순환형 MAPPO}

\subsection{반환값과 GAE}

팀 가치 예측을 \(V_t=V_\varphi(s_t)\), 종단 마스크를
\begin{equation}
  m_t=1-d_t
\end{equation}
로 둔다. 시간차 오차는
\begin{equation}
  \delta_t=
  r_t+\gamma m_tV_{t+1}-V_t
\end{equation}
이고 일반화 이점 추정치는
\begin{equation}
  \widehat A_t=
  \delta_t+
  \gamma\lambda m_t\widehat A_{t+1}.
  \label{eq:gae}
\end{equation}
반환값 표적은
\begin{equation}
  \widehat R_t=\widehat A_t+V_t
\end{equation}
다.

\subsection{PPO 확률비와 클립 목적}

로봇 \(i\)의 확률비는
\begin{equation}
  \rho_{t,i}(\theta)=
  \frac{
    \pi_\theta(a_i^t\mid o_i^t,h_i^{t-1},\Ftarget)
  }{
    \pi_{\theta_{\mathrm{old}}}
    (a_i^t\mid o_i^t,h_i^{t-1},\Ftarget)
  }
  \label{eq:probability-ratio}
\end{equation}
다. 공유 행위자의 클립 목적은
\begin{equation}
  \mathcal{J}_{\mathrm{clip}}(\theta)=
  \mathbb{E}_{t,i}\left[
    \min\left(
      \rho_{t,i}\widehat A_t,\,
      \clip(\rho_{t,i},1-\epsilon,1+\epsilon)\widehat A_t
    \right)
  \right].
  \label{eq:ppo-clip}
\end{equation}

\subsection{최종 손실}

정책 엔트로피를
\begin{equation}
  \mathcal{H}(\pi_\theta)=
  -\mathbb{E}_{t,i}
  \sum_a\pi_\theta(a\mid o_i^t)
  \log\pi_\theta(a\mid o_i^t)
\end{equation}
라 하면 행위자 손실은
\begin{equation}
  \mathcal{L}_{\mathrm{actor}}=
  -\mathcal{J}_{\mathrm{clip}}
  -c_H\mathcal{H}(\pi_\theta)
  +c_{\mathrm{aux}}\mathcal{L}_{\mathrm{aux}}.
  \label{eq:actor-loss}
\end{equation}

비평가 손실은
\begin{equation}
  \mathcal{L}_{\mathrm{critic}}=
  \mathbb{E}_t
  \left[
    \left(
      V_\varphi(s_t)-\widehat R_t
    \right)^2
  \right].
  \label{eq:critic-loss}
\end{equation}

순환 학습에서는 로봇별 시퀀스의 시간 순서를 유지한다. 에피소드 시작 지표 \(e_t\)에 대해
\begin{equation}
  h_i^{t-1}\leftarrow(1-e_t)h_i^{t-1}
\end{equation}
로 경계에서 은닉 상태를 초기화한다.

\section{커리큘럼과 도메인 무작위화}

커리큘럼 단계 \(c\in\{1,2,3,4\}\)에 따라 환경 파라미터 분포
\begin{equation}
  \zeta\sim p_c(\zeta)
\end{equation}
를 사용한다. \(\zeta\)에는
\begin{equation}
  \zeta=
  (w,\psi,A,\bm p_i^0,\theta_i^0,
   \mu,m,F_T,F_C,F_S,\sigma_s,\sigma_a)
\end{equation}
가 포함된다.

최근 \(W\)개 에피소드 성공률을
\begin{equation}
  \bar S_W=
  \frac{1}{W}
  \sum_{k=1}^{W}S_k
\end{equation}
라 하면 단계 전이는
\begin{equation}
  c\leftarrow
  \min(c+1,4)
  \quad\text{if}\quad
  \bar S_W\ge\tau_c
  \label{eq:curriculum-transition}
\end{equation}
다.

\section{모델 불변조건}

구현과 실험은 다음 불변조건을 유지해야 한다.

\begin{enumerate}[label=\arabic*.]
  \item 행위자 관측 \(o_i^t\)에는 전역 위치, 전체 지도, 전체 접촉 그래프와 현재 \(\Fcap\)이 없다.
  \item 모든 로봇은 동일한 \(\theta\)를 공유한다.
  \item 보상은 특정 형상 또는 토폴로지를 직접 평가하지 않는다.
  \item 성공은 \(\Ispan=1\)과 \(\Fcap\ge\Ftarget\)을 동시에 요구한다.
  \item 보조 표적은 훈련에만 쓰고 추론 입력으로 사용하지 않는다.
  \item 중앙 비평가의 특권 정보는 행위자 경로로 역으로 연결되지 않는다.
  \item 종단 용량은 로봇 행동을 동결한 별도 하중 시험으로 검증한다.
\end{enumerate}

\section{정식화의 범위}

본 모델은 창발적 집단 형태 형성과 정책 비교를 위한 연구 근사다. 식
\eqref{eq:normal-force}의 접촉, 식 \eqref{eq:anchor-failure}의 앵커,
식 \eqref{eq:fast-capacity} 및 \eqref{eq:terminal-capacity}의 구조 용량은
연속체 역학 또는 안전 인증 모델이 아니다. 향후 더 높은 충실도의 물리 모델로
대체할 때에도 식 \eqref{eq:decentralized-policy}, \eqref{eq:local-observation},
\eqref{eq:success}의 분산 실행 및 기능 성공 조건은 유지되어야 한다.

\end{document}
